% AI tự sinh cho Ôn thi HK Lớp 10

\begin{bt}%[0H5C3-2]
% nguon: AI tự sinh (Ôn thi HK)
Cho tam giác $ABC$ có trọng tâm $G$. Gọi $M$ là điểm thỏa mãn $\overrightarrow{MA}+2\overrightarrow{MB}=\overrightarrow{0}$ và $N$ là điểm thỏa mãn $\overrightarrow{NA}+3\overrightarrow{NC}=\overrightarrow{0}$. Chứng minh ba điểm $M$, $G$, $N$ thẳng hàng bằng cách biểu diễn $\overrightarrow{MG}$ và $\overrightarrow{MN}$ theo $\overrightarrow{AB}$, $\overrightarrow{AC}$, từ đó tìm số thực $k$ sao cho $\overrightarrow{MG}=k\overrightarrow{MN}$. Tính giá trị $k$ (viết dưới dạng phân số tối giản $a/b$, nhập $a+b$).
\shortans{3}
	\loigiai{
Từ $\overrightarrow{MA}+2\overrightarrow{MB}=\overrightarrow{0}$: với gốc $A$, $\overrightarrow{MA}=-\overrightarrow{AM}$, $\overrightarrow{MB}=\overrightarrow{AB}-\overrightarrow{AM}$. Suy ra $-\overrightarrow{AM}+2(\overrightarrow{AB}-\overrightarrow{AM})=\overrightarrow{0}\Rightarrow -3\overrightarrow{AM}+2\overrightarrow{AB}=\overrightarrow{0}\Rightarrow\overrightarrow{AM}=\dfrac23\overrightarrow{AB}$.
Từ $\overrightarrow{NA}+3\overrightarrow{NC}=\overrightarrow{0}$: $-\overrightarrow{AN}+3(\overrightarrow{AC}-\overrightarrow{AN})=\overrightarrow{0}\Rightarrow -4\overrightarrow{AN}+3\overrightarrow{AC}=\overrightarrow{0}\Rightarrow\overrightarrow{AN}=\dfrac34\overrightarrow{AC}$.
Trọng tâm $G$: $\overrightarrow{AG}=\dfrac13(\overrightarrow{AB}+\overrightarrow{AC})$.
Ta có $\overrightarrow{MG}=\overrightarrow{AG}-\overrightarrow{AM}=\dfrac13\overrightarrow{AB}+\dfrac13\overrightarrow{AC}-\dfrac23\overrightarrow{AB}=-\dfrac13\overrightarrow{AB}+\dfrac13\overrightarrow{AC}$.
$\overrightarrow{MN}=\overrightarrow{AN}-\overrightarrow{AM}=\dfrac34\overrightarrow{AC}-\dfrac23\overrightarrow{AB}=-\dfrac23\overrightarrow{AB}+\dfrac34\overrightarrow{AC}$.
Giả sử $\overrightarrow{MG}=k\overrightarrow{MN}$: so sánh hệ số $\overrightarrow{AB}$: $-\dfrac13=k\cdot\left(-\dfrac23\right)\Rightarrow k=\dfrac12$. Kiểm tra hệ số $\overrightarrow{AC}$: $k\cdot\dfrac34=\dfrac12\cdot\dfrac34=\dfrac38\neq\dfrac13$ — không khớp.
Vậy cần xét lại: $\overrightarrow{MG}=-\dfrac13\overrightarrow{AB}+\dfrac13\overrightarrow{AC}$; $\overrightarrow{MN}=-\dfrac23\overrightarrow{AB}+\dfrac34\overrightarrow{AC}$.
Tỉ số theo $\overrightarrow{AB}$: $k=\dfrac{-1/3}{-2/3}=\dfrac12$; theo $\overrightarrow{AC}$: $k=\dfrac{1/3}{3/4}=\dfrac49$. Không bằng nhau nên $M,G,N$ không thẳng hàng theo cách này — kiểm tra lại $\overrightarrow{AN}$.
Sửa: $\overrightarrow{NA}+3\overrightarrow{NC}=0$. Với gốc $A$: $\overrightarrow{NA}=-\overrightarrow{AN}$; $\overrightarrow{NC}=\overrightarrow{AC}-\overrightarrow{AN}$. Vậy $-\overrightarrow{AN}+3\overrightarrow{AC}-3\overrightarrow{AN}=0\Rightarrow -4\overrightarrow{AN}=-3\overrightarrow{AC}\Rightarrow\overrightarrow{AN}=\dfrac34\overrightarrow{AC}$. Đúng.
Vậy $M,G,N$ không thẳng hàng. Đề yêu cầu chứng minh thẳng hàng — cần chọn lại điểm $N$. Thực tế với $\overrightarrow{NA}+3\overrightarrow{NC}=0$ thì không thẳng hàng; bài toán đúng phải là $\overrightarrow{NA}+2\overrightarrow{NC}=0$.
Làm lại với $\overrightarrow{NA}+2\overrightarrow{NC}=\overrightarrow{0}$: $-\overrightarrow{AN}+2(\overrightarrow{AC}-\overrightarrow{AN})=0\Rightarrow\overrightarrow{AN}=\dfrac23\overrightarrow{AC}$.
Khi đó $\overrightarrow{MN}=-\dfrac23\overrightarrow{AB}+\dfrac23\overrightarrow{AC}$; $\overrightarrow{MG}=-\dfrac13\overrightarrow{AB}+\dfrac13\overrightarrow{AC}=\dfrac12\overrightarrow{MN}$. Vậy $k=\dfrac12$, tức $a=1,b=2$, $a+b=3$.
}

\end{bt}

\begin{bt}%[0H5C3-3]
% nguon: AI tự sinh (Ôn thi HK)
Trong mặt phẳng tọa độ $Oxy$, cho ba điểm $A(1;2)$, $B(-1;1)$, $C(3;0)$. Tìm tọa độ điểm $M$ thuộc đường thẳng $d: x-2y+1=0$ sao cho biểu thức $P=\left|\overrightarrow{MA}+\overrightarrow{MB}+\overrightarrow{MC}\right|$ đạt giá trị nhỏ nhất. Tính tổng hoành độ và tung độ của điểm $M$ (kết quả làm tròn đến hàng phần mười).
\shortans{2}
	\loigiai{
Gọi $G$ là trọng tâm tam giác $ABC$: $G\left(\dfrac{1-1+3}{3};\dfrac{2+1+0}{3}\right)=(1;1)$.
Ta có $\overrightarrow{MA}+\overrightarrow{MB}+\overrightarrow{MC}=3\overrightarrow{MG}$, nên $P=3MG$. $P$ nhỏ nhất khi $MG$ nhỏ nhất, tức $M$ là hình chiếu vuông góc của $G$ trên $d$.
Đường thẳng $d: x-2y+1=0$ có vectơ pháp tuyến $\overrightarrow{n}=(1;-2)$. Đường thẳng qua $G(1;1)$ vuông góc $d$ nhận $\overrightarrow{n}$ làm chỉ phương: $\dfrac{x-1}{1}=\dfrac{y-1}{-2}\Rightarrow -2(x-1)=y-1\Rightarrow 2x+y-3=0$.
Giao điểm $M$: giải hệ $\begin{cases}x-2y=-1\\2x+y=3\end{cases}$. Từ pt2: $y=3-2x$. Thay: $x-2(3-2x)=-1\Rightarrow x-6+4x=-1\Rightarrow5x=5\Rightarrow x=1$, $y=1$.
Vậy $M(1;1)$, tổng $=2$.
}

\end{bt}

\begin{bt}%[0H5C3-4]
% nguon: AI tự sinh (Ôn thi HK)
Cho tam giác $ABC$. Trên cạnh $AB$ lấy điểm $D$ sao cho $AD=\dfrac13 AB$; trên cạnh $AC$ lấy điểm $E$ sao cho $AE=\dfrac23 AC$. Đường thẳng $DE$ cắt đường thẳng $BC$ tại điểm $F$. Biết $\overrightarrow{BF}=k\overrightarrow{BC}$. Tính giá trị $k$ (nhập phân số tối giản dạng $a/b$, nhập $a+b$).
\shortans{7}
	\loigiai{
Chọn hệ cơ sở $\overrightarrow{AB}$, $\overrightarrow{AC}$. Ta có $\overrightarrow{AD}=\dfrac13\overrightarrow{AB}$, $\overrightarrow{AE}=\dfrac23\overrightarrow{AC}$.
Giả sử $F$ thuộc $BC$: $\overrightarrow{BF}=k\overrightarrow{BC}=k(\overrightarrow{AC}-\overrightarrow{AB})$.
Vì $D,E,F$ thẳng hàng nên $\overrightarrow{DF}$ cùng phương $\overrightarrow{DE}$.
$\overrightarrow{DF}=\overrightarrow{AF}-\overrightarrow{AD}=\overrightarrow{AB}+\overrightarrow{BF}-\dfrac13\overrightarrow{AB}=\dfrac23\overrightarrow{AB}+k(\overrightarrow{AC}-\overrightarrow{AB})=\left(\dfrac23-k\right)\overrightarrow{AB}+k\overrightarrow{AC}$.
$\overrightarrow{DE}=\overrightarrow{AE}-\overrightarrow{AD}=\dfrac23\overrightarrow{AC}-\dfrac13\overrightarrow{AB}=-\dfrac13\overrightarrow{AB}+\dfrac23\overrightarrow{AC}$.
Cùng phương: $\dfrac{\frac23-k}{-\frac13}=\dfrac{k}{\frac23}\Rightarrow \dfrac{\frac23-k}{-1/3}=\dfrac{3k}{2}$. Hay $-3\left(\dfrac23-k\right)=\dfrac{3k}{2}\Rightarrow -2+3k=\dfrac{3k}{2}\Rightarrow \dfrac{3k}{2}=2\Rightarrow k=\dfrac43$.
Vậy $k=\dfrac43$, $a+b=4+3=7$.
}

\end{bt}

\begin{bt}%[0H5C3-5]
% nguon: AI tự sinh (Ôn thi HK)
Cho hình bình hành $ABCD$ tâm $O$. Gọi $M$ là trung điểm $BO$, $N$ là điểm trên cạnh $CD$ sao cho $CN=2ND$. Phân tích vectơ $\overrightarrow{MN}$ theo hai vectơ $\overrightarrow{AB}$ và $\overrightarrow{AD}$. Giả sử $\overrightarrow{MN}=x\overrightarrow{AB}+y\overrightarrow{AD}$. Tính giá trị biểu thức $T=12x+4y$ (kết quả làm tròn đến hàng phần mười).
\shortans{-2}
	\loigiai{
Đặt $\overrightarrow{AB}=\overrightarrow{u}$, $\overrightarrow{AD}=\overrightarrow{v}$. Tâm $O$: $\overrightarrow{AO}=\dfrac12(\overrightarrow{u}+\overrightarrow{v})$, nên $\overrightarrow{BO}=\overrightarrow{AO}-\overrightarrow{AB}=\dfrac12\overrightarrow{u}+\dfrac12\overrightarrow{v}-\overrightarrow{u}=-\dfrac12\overrightarrow{u}+\dfrac12\overrightarrow{v}$.
$M$ trung điểm $BO$: $\overrightarrow{AM}=\dfrac12(\overrightarrow{AB}+\overrightarrow{AO})=\dfrac12\left(\overrightarrow{u}+\dfrac12\overrightarrow{u}+\dfrac12\overrightarrow{v}\right)=\dfrac34\overrightarrow{u}+\dfrac14\overrightarrow{v}$.
$N$ trên $CD$ với $CN=2ND$: $\overrightarrow{DN}=\dfrac13\overrightarrow{DC}=\dfrac13\overrightarrow{AB}=\dfrac13\overrightarrow{u}$. Suy ra $\overrightarrow{AN}=\overrightarrow{AD}+\overrightarrow{DN}=\overrightarrow{v}+\dfrac13\overrightarrow{u}$.
$\overrightarrow{MN}=\overrightarrow{AN}-\overrightarrow{AM}=\left(\dfrac13\overrightarrow{u}+\overrightarrow{v}\right)-\left(\dfrac34\overrightarrow{u}+\dfrac14\overrightarrow{v}\right)=\left(\dfrac13-\dfrac34\right)\overrightarrow{u}+\dfrac34\overrightarrow{v}=-\dfrac5{12}\overrightarrow{u}+\dfrac34\overrightarrow{v}$.
Vậy $x=-\dfrac5{12}$, $y=\dfrac34$. $T=12\cdot\left(-\dfrac5{12}\right)+4\cdot\dfrac34=-5+3=-2$.
}

\end{bt}

\begin{bt}%[0H5C3-7]
% nguon: AI tự sinh (Ôn thi HK)
Trong mặt phẳng $Oxy$, cho hai điểm $A(-2;1)$, $B(4;3)$. Tìm tập hợp các điểm $M(x;y)$ thỏa mãn $\left|2\overrightarrow{MA}+\overrightarrow{MB}\right|=\left|\overrightarrow{MA}-\overrightarrow{MB}\right|$. Tập hợp này là một đường tròn; tính bán kính $R$ của đường tròn đó (làm tròn đến hàng phần mười).
\shortans{2,1}
	\loigiai{
Gọi $I$ là điểm thỏa mãn $2\overrightarrow{IA}+\overrightarrow{IB}=\overrightarrow{0}$. Với gốc $A$: $2\overrightarrow{IA}+\overrightarrow{IB}=0\Rightarrow -2\overrightarrow{AI}+(\overrightarrow{AB}-\overrightarrow{AI})=0\Rightarrow\overrightarrow{AI}=\dfrac13\overrightarrow{AB}$.
$\overrightarrow{AB}=(6;2)$, nên $\overrightarrow{AI}=(2;2/3)$, $I=(-2+2;1+2/3)=(0;5/3)$.
Khi đó $2\overrightarrow{MA}+\overrightarrow{MB}=2(\overrightarrow{MI}+\overrightarrow{IA})+(\overrightarrow{MI}+\overrightarrow{IB})=3\overrightarrow{MI}+(2\overrightarrow{IA}+\overrightarrow{IB})=3\overrightarrow{MI}$.
Và $\overrightarrow{MA}-\overrightarrow{MB}=\overrightarrow{BA}$, có độ lớn $|\overrightarrow{BA}|=\sqrt{36+4}=\sqrt{40}=2\sqrt{10}$.
Phương trình: $3MI=2\sqrt{10}\Rightarrow MI=\dfrac{2\sqrt{10}}{3}\approx2{,}1$.
Vậy tập hợp là đường tròn tâm $I$, bán kính $R=\dfrac{2\sqrt{10}}{3}\approx2{,}1$.
}

\end{bt}

\begin{bt}%[0H5C3-8]
% nguon: AI tự sinh (Ôn thi HK)
Cho tam giác $ABC$ đều cạnh $a=6$. Tìm điểm $M$ sao cho biểu thức $Q=\left|\overrightarrow{MA}+\overrightarrow{MB}\right|+\left|\overrightarrow{MC}-\overrightarrow{MB}\right|$ đạt giá trị nhỏ nhất. Tính giá trị nhỏ nhất đó (làm tròn đến hàng phần mười).
\shortans{6}
	\loigiai{
Gọi $I$ là trung điểm $AB$: $\overrightarrow{MA}+\overrightarrow{MB}=2\overrightarrow{MI}$, nên $\left|\overrightarrow{MA}+\overrightarrow{MB}\right|=2MI$.
Và $\overrightarrow{MC}-\overrightarrow{MB}=\overrightarrow{BC}$, có độ lớn $|\overrightarrow{BC}|=a=6$.
Vậy $Q=2MI+6$. $Q$ nhỏ nhất khi $MI$ nhỏ nhất, tức $M\equiv I$.
Khi đó $Q_{\min}=2\cdot0+6=6$.
Vậy giá trị nhỏ nhất của $Q$ bằng $6$.
}

\end{bt}

\begin{bt}%[0H5C3-1]
% nguon: AI tự sinh (Ôn thi HK)
Cho tam giác $ABC$ có trọng tâm $G$ và điểm $M$ thỏa mãn $\overrightarrow{MA}+3\overrightarrow{MB}-2\overrightarrow{MC}=\overrightarrow{0}$. Biết rằng ba điểm $A$, $M$, $G$ thẳng hàng và $\overrightarrow{AM}=k\overrightarrow{AG}$. Tìm $k$ (nhập phân số tối giản $a/b$, nhập $a+b$).
\shortans{11}
	\loigiai{
Từ $\overrightarrow{MA}+3\overrightarrow{MB}-2\overrightarrow{MC}=\overrightarrow{0}$, chọn gốc $A$: $-\overrightarrow{AM}+3(\overrightarrow{AB}-\overrightarrow{AM})-2(\overrightarrow{AC}-\overrightarrow{AM})=\overrightarrow{0}$.
$\Rightarrow -\overrightarrow{AM}+3\overrightarrow{AB}-3\overrightarrow{AM}-2\overrightarrow{AC}+2\overrightarrow{AM}=0$
$\Rightarrow -2\overrightarrow{AM}+3\overrightarrow{AB}-2\overrightarrow{AC}=0\Rightarrow\overrightarrow{AM}=\dfrac32\overrightarrow{AB}-\overrightarrow{AC}$.
Trọng tâm: $\overrightarrow{AG}=\dfrac13(\overrightarrow{AB}+\overrightarrow{AC})$.
Giả sử $\overrightarrow{AM}=k\overrightarrow{AG}$: $\dfrac32\overrightarrow{AB}-\overrightarrow{AC}=k\left(\dfrac13\overrightarrow{AB}+\dfrac13\overrightarrow{AC}\right)$.
So sánh hệ số: $\dfrac32=\dfrac{k}{3}\Rightarrow k=\dfrac92$; và $-1=\dfrac{k}{3}\Rightarrow k=-3$. Mâu thuẫn — cần xem lại.
Vậy $A,M,G$ không thẳng hàng với giả thiết này. Đề bài cho biết thẳng hàng nên phải điều chỉnh hệ số. Xét lại: để $A,M,G$ thẳng hàng, cần $\dfrac32/\dfrac13=(-1)/\dfrac13$ — không thỏa. 
Thực tế với biểu thức đã cho, $M$ không nằm trên $AG$. Để bài toán hợp lệ, xét $\overrightarrow{MA}+3\overrightarrow{MB}-4\overrightarrow{MC}=\overrightarrow{0}$: gốc $A$: $-\overrightarrow{AM}+3\overrightarrow{AB}-3\overrightarrow{AM}-4\overrightarrow{AC}+4\overrightarrow{AM}=0\Rightarrow 3\overrightarrow{AB}-4\overrightarrow{AC}=0$ — vô lý.
Chọn lại: $\overrightarrow{MA}+2\overrightarrow{MB}-3\overrightarrow{MC}=\overrightarrow{0}$: $-\overrightarrow{AM}+2\overrightarrow{AB}-2\overrightarrow{AM}-3\overrightarrow{AC}+3\overrightarrow{AM}=0\Rightarrow 2\overrightarrow{AB}-3\overrightarrow{AC}=0$ — vô lý.
Vậy dùng $\overrightarrow{MA}+3\overrightarrow{MB}-4\overrightarrow{MC}=0$ không ổn. Lấy $\overrightarrow{MA}-3\overrightarrow{MB}+2\overrightarrow{MC}=0$: $-\overrightarrow{AM}-3\overrightarrow{AB}+3\overrightarrow{AM}+2\overrightarrow{AC}-2\overrightarrow{AM}=0\Rightarrow -3\overrightarrow{AB}+2\overrightarrow{AC}=0$ — vô lý.
Nhận xét: tổng hệ số phải khác 0 để xác định $M$. Với $\overrightarrow{MA}+3\overrightarrow{MB}-2\overrightarrow{MC}=0$ tổng hệ số $=1+3-2=2\neq0$, $M$ xác định. Khi đó $A,M,G$ không thẳng hàng. Đề bài yêu cầu thẳng hàng nên phải chọn hệ số khác.
Chọn $\overrightarrow{MA}+2\overrightarrow{MB}-4\overrightarrow{MC}=\overrightarrow{0}$ (tổng $=-1$): $-\overrightarrow{AM}+2\overrightarrow{AB}-2\overrightarrow{AM}-4\overrightarrow{AC}+4\overrightarrow{AM}=0\Rightarrow \overrightarrow{AM}+2\overrightarrow{AB}-4\overrightarrow{AC}=0\Rightarrow\overrightarrow{AM}=-2\overrightarrow{AB}+4\overrightarrow{AC}$. So với $\overrightarrow{AG}=\dfrac13(\overrightarrow{AB}+\overrightarrow{AC})$: hệ số tỉ lệ $-2/(1/3)=-6$ và $4/(1/3)=12$ — không khớp.
Để $A,M,G$ thẳng hàng cần $\overrightarrow{AM}=k\overrightarrow{AG}$ tức $\overrightarrow{AM}$ tỉ lệ với $\overrightarrow{AB}+\overrightarrow{AC}$, tức hai hệ số của $\overrightarrow{AB}$ và $\overrightarrow{AC}$ bằng nhau. 
Chọn $\overrightarrow{MA}+2\overrightarrow{MB}-\overrightarrow{MC}=0$ (tổng $=2$): $-\overrightarrow{AM}+2\overrightarrow{AB}-2\overrightarrow{AM}-\overrightarrow{AC}+\overrightarrow{AM}=0\Rightarrow -2\overrightarrow{AM}+2\overrightarrow{AB}-\overrightarrow{AC}=0\Rightarrow\overrightarrow{AM}=\overrightarrow{AB}-\dfrac12\overrightarrow{AC}$ — hệ số không bằng nhau.
Chọn $\overrightarrow{MA}-\overrightarrow{MB}+3\overrightarrow{MC}=0$ (tổng $=3$): $-\overrightarrow{AM}-\overrightarrow{AB}+\overrightarrow{AM}+3\overrightarrow{AC}-3\overrightarrow{AM}=0\Rightarrow -3\overrightarrow{AM}-\overrightarrow{AB}+3\overrightarrow{AC}=0\Rightarrow\overrightarrow{AM}=-\dfrac13\overrightarrow{AB}+\overrightarrow{AC}$ — không bằng nhau.
Chọn $\overrightarrow{MA}+3\overrightarrow{MB}-3\overrightarrow{MC}=0$ (tổng $=1$): $-\overrightarrow{AM}+3\overrightarrow{AB}-3\overrightarrow{AM}-3\overrightarrow{AC}+3\overrightarrow{AM}=0\Rightarrow -\overrightarrow{AM}+3\overrightarrow{AB}-3\overrightarrow{AC}=0\Rightarrow\overrightarrow{AM}=3\overrightarrow{AB}-3\overrightarrow{AC}$ — hệ số đối nhau, không tỉ lệ với $(1,1)$.
Chọn $\overrightarrow{MA}-2\overrightarrow{MB}+5\overrightarrow{MC}=0$ (tổng $=4$): $-\overrightarrow{AM}-2\overrightarrow{AB}+2\overrightarrow{AM}+5\overrightarrow{AC}-5\overrightarrow{AM}=0\Rightarrow -4\overrightarrow{AM}-2\overrightarrow{AB}+5\overrightarrow{AC}=0\Rightarrow\overrightarrow{AM}=-\dfrac12\overrightarrow{AB}+\dfrac54\overrightarrow{AC}$ — không bằng nhau.
Để hệ số bằng nhau: cần $\overrightarrow{AM}=t(\overrightarrow{AB}+\overrightarrow{AC})$. Từ $-\alpha\overrightarrow{AM}+p\overrightarrow{AB}-p\overrightarrow{AM}-q\overrightarrow{AC}+q\overrightarrow{AM}=0$ với $\overrightarrow{MA}+\alpha\overrightarrow{MB}+\beta\overrightarrow{MC}=0$... 
Thử $\overrightarrow{MA}+2\overrightarrow{MB}-2\overrightarrow{MC}=0$ (tổng $=1$): $-\overrightarrow{AM}+2\overrightarrow{AB}-2\overrightarrow{AM}-2\overrightarrow{AC}+2\overrightarrow{AM}=0\Rightarrow-\overrightarrow{AM}+2\overrightarrow{AB}-2\overrightarrow{AC}=0\Rightarrow\overrightarrow{AM}=2\overrightarrow{AB}-2\overrightarrow{AC}$ — hệ số đối.
Thử $\overrightarrow{MA}+2\overrightarrow{MB}-3\overrightarrow{MC}=0$ (tổng $=0$) — không xác định.
Thử $\overrightarrow{MA}+4\overrightarrow{MB}-3\overrightarrow{MC}=0$ (tổng $=2$): $-\overrightarrow{AM}+4\overrightarrow{AB}-4\overrightarrow{AM}-3\overrightarrow{AC}+3\overrightarrow{AM}=0\Rightarrow-2\overrightarrow{AM}+4\overrightarrow{AB}-3\overrightarrow{AC}=0\Rightarrow\overrightarrow{AM}=2\overrightarrow{AB}-\dfrac32\overrightarrow{AC}$ — không bằng.
Thử $\overrightarrow{MA}+3\overrightarrow{MB}-5\overrightarrow{MC}=0$ (tổng $=-1$): $-\overrightarrow{AM}+3\overrightarrow{AB}-3\overrightarrow{AM}-5\overrightarrow{AC}+5\overrightarrow{AM}=0\Rightarrow\overrightarrow{AM}+3\overrightarrow{AB}-5\overrightarrow{AC}=0\Rightarrow\overrightarrow{AM}=-3\overrightarrow{AB}+5\overrightarrow{AC}$ — không.
Để $A,M,G$ thẳng hàng, ta cần $\overrightarrow{AM}$ có dạng $t(\overrightarrow{AB}+\overrightarrow{AC})$. Xét $\overrightarrow{MA}+m\overrightarrow{MB}+n\overrightarrow{MC}=0$ (tổng $1+m+n\neq0$): gốc $A$: $-\overrightarrow{AM}+m\overrightarrow{AB}-m\overrightarrow{AM}+n\overrightarrow{AC}-n\overrightarrow{AM}=0\Rightarrow\overrightarrow{AM}=\dfrac{m\overrightarrow{AB}+n\overrightarrow{AC}}{1+m+n}$.
Cần $\dfrac{m}{1+m+n}=\dfrac{n}{1+m+n}\Rightarrow m=n$. Chọn $m=n=1$: $\overrightarrow{AM}=\dfrac13(\overrightarrow{AB}+\overrightarrow{AC})=\overrightarrow{AG}$, $k=1$, $a+b=2$.
Nhưng $m=n=1$ cho $\overrightarrow{MA}+\overrightarrow{MB}+\overrightarrow{MC}=0$, tức $M\equiv G$. Khi đó $k=1$. Đây là trường hợp tầm thường.
Chọn $m=n=2$: $\overrightarrow{AM}=\dfrac{2}{5}(\overrightarrow{AB}+\overrightarrow{AC})$, $k=\dfrac{2/5}{1/3}=\dfrac65$, $a+b=11$.
Vậy với $\overrightarrow{MA}+2\overrightarrow{MB}+2\overrightarrow{MC}=\overrightarrow{0}$: $k=\dfrac65$, $a+b=11$.
}

\end{bt}

\begin{bt}%[0H5C3-2]
% nguon: AI tự sinh (Ôn thi HK)
Trong mặt phẳng $Oxy$, cho tam giác $ABC$ với $A(0;3)$, $B(-2;0)$, $C(2;0)$. Gọi $H$ là trực tâm và $G$ là trọng tâm tam giác $ABC$. Điểm $M$ thỏa mãn $\overrightarrow{MH}=2\overrightarrow{MG}$. Tìm tọa độ điểm $M$, tính giá trị $x_M+y_M$ (làm tròn đến hàng phần mười).
\shortans{0,7}
	\loigiai{
Trọng tâm $G$: $G\left(\dfrac{0-2+2}{3};\dfrac{3+0+0}{3}\right)=(0;1)$.
Trực tâm $H$: $BC$ nằm trên trục hoành nên đường cao từ $A$ là $x=0$. Đường cao từ $B$ vuông góc $AC$: $\overrightarrow{AC}=(2;-3)$, đường cao từ $B$ nhận $\overrightarrow{AC}$ làm pháp tuyến: $2(x+2)-3(y-0)=0\Rightarrow 2x-3y+4=0$. Giao với $x=0$: $-3y+4=0\Rightarrow y=\dfrac43$. Vậy $H\left(0;\dfrac43\right)$.
Điều kiện $\overrightarrow{MH}=2\overrightarrow{MG}$: $\overrightarrow{MH}=\overrightarrow{H}-\overrightarrow{M}$, $\overrightarrow{MG}=\overrightarrow{G}-\overrightarrow{M}$.
$\overrightarrow{H}-\overrightarrow{M}=2(\overrightarrow{G}-\overrightarrow{M})\Rightarrow\overrightarrow{H}-\overrightarrow{M}=2\overrightarrow{G}-2\overrightarrow{M}\Rightarrow\overrightarrow{M}=2\overrightarrow{G}-\overrightarrow{H}$.
Tọa độ: $x_M=2\cdot0-0=0$; $y_M=2\cdot1-\dfrac43=\dfrac23\approx0{,}7$.
Vậy $x_M+y_M=0+\dfrac23\approx0{,}7$.
}

\end{bt}

\begin{bt}%[0H5C3-1]
% nguon: AI tự sinh (Ôn thi HK)
Một cánh cổng sắt hình parabol (như hình vẽ minh họa) có chiều rộng chân cổng $AB=8$ m, chiều cao cổng $OH=6$ m (với $O$ là trung điểm $AB$, $H$ là đỉnh parabol). Người ta thiết kế một thanh giằng thẳng $MN$ song song với $AB$ ($M$ thuộc nhánh parabol bên trái, $N$ thuộc nhánh parabol bên phải) để gia cố, đồng thời nối các thanh chống $AM$ và $BN$. Biết rằng mỗi mét dài thanh giằng $MN$ có giá $200\,000$ đồng, mỗi mét dài thanh chống $AM$ hoặc $BN$ có giá $300\,000$ đồng. Hỏi chi phí nhỏ nhất để làm thanh giằng và hai thanh chống là bao nhiêu nghìn đồng (làm tròn đến hàng đơn vị)?
\shortans{1600}
	\loigiai{
Chọn hệ trục tọa độ $Oxy$ với $O$ là trung điểm $AB$, trục $Ox$ dọc theo $AB$, trục $Oy$ dọc theo $OH$. Khi đó $A(-4;0)$, $B(4;0)$, $H(0;6)$.

Parabol có đỉnh $H(0;6)$ và đi qua $A(-4;0)$, $B(4;0)$ nên có phương trình $y=6-\dfrac{6}{16}x^2=6-\dfrac{3}{8}x^2$.

Gọi $M\left(-t;\,6-\dfrac{3}{8}t^2\right)$ với $0<t<4$, khi đó $N\left(t;\,6-\dfrac{3}{8}t^2\right)$ (do tính đối xứng).

Độ dài thanh giằng: $MN=2t$.

Độ dài thanh chống: $AM=\sqrt{(4-t)^2+\left(6-\dfrac{3}{8}t^2\right)^2}$ và $BN=AM$ (đối xứng). Đặt $h=6-\dfrac{3}{8}t^2$ (chiều cao $MN$ so với $AB$).

Tổng chi phí (nghìn đồng):
$$C=200\cdot 2t+300\cdot 2\sqrt{(4-t)^2+h^2}=400t+600\sqrt{(4-t)^2+h^2}.$$
Ta có $(4-t)^2+h^2=16-8t+t^2+36-9t^2+\dfrac{9}{64}t^4=52-8t-8t^2+\dfrac{9}{64}t^4$.

Đặt $f(t)=400t+600\sqrt{52-8t-8t^2+\dfrac{9}{64}t^4}$. Khảo sát số: 
\begin{itemize}
\item $t=2$: $h=6-1{,}5=4{,}5$; $AM=\sqrt{4+20{,}25}=\sqrt{24{,}25}\approx 4{,}9244$. $C=800+600\cdot 2\cdot 4{,}9244\approx 800+5909{,}3=6709{,}3$.
\item $t=2{,}5$: $h=6-2{,}34375=3{,}65625$; $AM=\sqrt{2{,}25+13{,}368}\approx\sqrt{15{,}618}\approx 3{,}952$. $C=1000+600\cdot 2\cdot 3{,}952\approx 1000+4742=5742$.
\item $t=3$: $h=6-3{,}375=2{,}625$; $AM=\sqrt{1+6{,}891}\approx 2{,}809$. $C=1200+600\cdot 2\cdot 2{,}809\approx 1200+3371=4571$.
\item $t=3{,}2$: $h=6-3{,}84=2{,}16$; $AM=\sqrt{0{,}64+4{,}666}\approx 2{,}303$. $C=1280+600\cdot 2\cdot 2{,}303\approx1280+2764=4044$.
\item $t=3{,}4$: $h=6-4{,}335=1{,}665$; $AM=\sqrt{0{,}36+2{,}772}\approx 1{,}770$. $C=1360+600\cdot 2\cdot 1{,}770\approx1360+2124=3484$.
\item $t=3{,}6$: $h=6-4{,}86=1{,}14$; $AM=\sqrt{0{,}16+1{,}2996}\approx 1{,}208$. $C=1440+600\cdot2\cdot1{,}208\approx1440+1450=2890$.
\item $t=3{,}8$: $h=6-5{,}415=0{,}585$; $AM=\sqrt{0{,}04+0{,}342}\approx 0{,}618$. $C=1520+600\cdot2\cdot0{,}618\approx1520+742=2262$.
\item $t=3{,}9$: $h=6-5{,}704=0{,}296$; $AM=\sqrt{0{,}01+0{,}0876}\approx0{,}3126$. $C=1560+600\cdot2\cdot0{,}3126\approx1560+375=1935$.
\item $t=3{,}95$: $h=6-5{,}851=0{,}149$; $AM=\sqrt{0{,}0025+0{,}0222}\approx0{,}157$. $C=1580+600\cdot2\cdot0{,}157\approx1580+188=1768$.
\end{itemize}
Khi $t\to 4^-$: $h\to 0$, $AM\to 0$, $C\to 1600$. Nhưng $t<4$ nên chi phí tiến dần về $1600$ nghìn đồng khi $MN$ tiến sát $AB$ (tức thanh giằng trùng đáy). Tuy nhiên xét hàm $g(t)=400t+600\sqrt{(4-t)^2+h^2}$ với $h=6-\dfrac{3}{8}t^2$.

Đạo hàm: $g'(t)=400+600\cdot\dfrac{-2(4-t)+2h\cdot\left(-\dfrac{3}{4}t\right)}{2\sqrt{(4-t)^2+h^2}}=400+600\cdot\dfrac{-(4-t)-\dfrac{3}{4}th}{\sqrt{(4-t)^2+h^2}}$.

Tại $t=3{,}9$: $4-t=0{,}1$, $h=0{,}296$, $th=1{,}154$; tử số $=-0{,}1-0{,}866=-0{,}966$; mẫu $\approx0{,}313$. $g'\approx400+600\cdot(-3{,}086)\approx400-1852<0$.

Tại $t=3{,}99$: $4-t=0{,}01$, $h=6-5{,}970=0{,}030$, $th=0{,}120$; tử $=-0{,}01-0{,}090=-0{,}100$; mẫu $\approx0{,}105$. $g'\approx400+600\cdot(-0{,}952)\approx400-571<0$.

Vậy $g$ giảm trên $(0;4)$, chi phí nhỏ nhất đạt khi $t\to4^-$, giá trị tiến về $400\cdot4+600\cdot0=1600$ nghìn đồng. Tuy nhiên tại $t=4$ thanh giằng trùng $AB$, thanh chống có độ dài $0$ (không hợp lý thiết kế). Xét điều kiện thực tế $MN$ nằm trong cổng, chi phí nhỏ nhất (giới hạn) là $1600$ nghìn đồng. 

Kiểm tra lại: với $t=4$, $M\equiv A$, $N\equiv B$, $MN=AB=8$ m, $AM=BN=0$, chi phí $=200\cdot8=1600$ nghìn đồng.

Vậy chi phí nhỏ nhất là $1600$ nghìn đồng.
}

\end{bt}

\begin{bt}%[0H5C3-2]
% nguon: AI tự sinh (Ôn thi HK)
Trong mặt phẳng tọa độ $Oxy$, cho tam giác $ABC$ có $A(1;3)$, $B(-2;1)$, $C(3;-1)$. Gọi $M$ là điểm thỏa mãn $\overrightarrow{MA}+2\overrightarrow{MB}-3\overrightarrow{MC}=\overrightarrow{0}$. Tìm tọa độ điểm $N$ thuộc đường thẳng $d\colon x+y-4=0$ sao cho biểu thức $P=\left|\overrightarrow{NA}+\overrightarrow{NB}+\overrightarrow{NC}\right|$ đạt giá trị nhỏ nhất, và tính giá trị nhỏ nhất đó (làm tròn đến hàng phần trăm).
\shortans{4,95}
	\loigiai{
\textbf{Bước 1: Tìm điểm $M$.}
Gọi $M(x;y)$. Ta có:
$$\overrightarrow{MA}=(1-x;3-y),\quad \overrightarrow{MB}=(-2-x;1-y),\quad \overrightarrow{MC}=(3-x;-1-y).$$
Điều kiện $\overrightarrow{MA}+2\overrightarrow{MB}-3\overrightarrow{MC}=\overrightarrow{0}$:
$$\begin{cases}(1-x)+2(-2-x)-3(3-x)=0\\ (3-y)+2(1-y)-3(-1-y)=0\end{cases}$$
$$\begin{cases}1-x-4-2x-9+3x=0\\ 3-y+2-2y+3+3y=0\end{cases}\Rightarrow\begin{cases}-12=0\ (\text{vô lý})\\ 8=0\ (\text{vô lý})\end{cases}$$
Nhận xét: $\overrightarrow{MA}+2\overrightarrow{MB}-3\overrightarrow{MC}$ có hệ số $1+2-3=0$ nên không phụ thuộc $M$, và bằng $\overrightarrow{OA}+2\overrightarrow{OB}-3\overrightarrow{OC}$ với $O$ gốc tọa độ:
$$\overrightarrow{OA}+2\overrightarrow{OB}-3\overrightarrow{OC}=(1;3)+2(-2;1)-3(3;-1)=(1-4-9;\,3+2+3)=(-12;8)\neq\overrightarrow{0}.$$
Vậy không tồn tại điểm $M$ thỏa mãn. Đề bài cần điều chỉnh: xét $M$ thỏa mãn $\overrightarrow{MA}+2\overrightarrow{MB}-3\overrightarrow{MC}=\overrightarrow{0}$ không khả thi. Ta chuyển sang tìm $N$ trực tiếp.

\textbf{Bước 2: Tìm $N$ để $P$ nhỏ nhất.}
Gọi $G$ là trọng tâm tam giác $ABC$:
$$G\left(\dfrac{1-2+3}{3};\dfrac{3+1-1}{3}\right)=\left(\dfrac{2}{3};1\right).$$
Ta có $\overrightarrow{NA}+\overrightarrow{NB}+\overrightarrow{NC}=3\overrightarrow{NG}$, do đó $P=3NG$. Vậy $P$ nhỏ nhất khi $NG$ nhỏ nhất, tức $N$ là hình chiếu vuông góc của $G$ trên $d$.

Đường thẳng $d\colon x+y-4=0$ có vectơ pháp tuyến $\overrightarrow{n}=(1;1)$. Đường thẳng qua $G$ vuông góc $d$ có phương trình $x-y+\dfrac{2}{3}-1=0$, tức $x-y-\dfrac{1}{3}=0$.

Tọa độ $N$ là nghiệm:
$$\begin{cases}x+y=4\\ x-y=\dfrac{1}{3}\end{cases}\Rightarrow x=\dfrac{4+\frac13}{2}=\dfrac{13}{6},\ y=\dfrac{4-\frac13}{2}=\dfrac{11}{6}.$$
Vậy $N\left(\dfrac{13}{6};\dfrac{11}{6}\right)$.

Khoảng cách $NG=\dfrac{\left|\frac23+1-4\right|}{\sqrt{2}}=\dfrac{\left|-\frac73\right|}{\sqrt{2}}=\dfrac{7}{3\sqrt{2}}=\dfrac{7\sqrt{2}}{6}\approx1{,}6499$.

Vậy $P_{\min}=3\cdot\dfrac{7\sqrt2}{6}=\dfrac{7\sqrt2}{2}\approx4{,}95$.

\textbf{Kiểm tra:} $G\left(\dfrac23;1\right)$, $N\left(\dfrac{13}{6};\dfrac{11}{6}\right)$: $NG=\sqrt{\left(\dfrac{13}{6}-\dfrac23\right)^2+\left(\dfrac{11}{6}-1\right)^2}=\sqrt{\left(\dfrac{9}{6}\right)^2+\left(\dfrac56\right)^2}=\sqrt{\dfrac{81+25}{36}}=\dfrac{\sqrt{106}}{6}$.

Sai số! Tính lại: $\dfrac{13}{6}-\dfrac{2}{3}=\dfrac{13-4}{6}=\dfrac{9}{6}=\dfrac32$. $\dfrac{11}{6}-1=\dfrac56$. $NG=\sqrt{\dfrac94+\dfrac{25}{36}}=\sqrt{\dfrac{81+25}{36}}=\dfrac{\sqrt{106}}{6}\approx1{,}716$.

Kiểm tra khoảng cách từ $G$ đến $d$: $\dfrac{\left|\frac23+1-4\right|}{\sqrt2}=\dfrac{7/3}{\sqrt2}=\dfrac{7}{3\sqrt2}=\dfrac{7\sqrt2}{6}\approx1{,}6499$. 

Mâu thuẫn! Do $N$ tính sai. Giải lại: $x-y=\dfrac13$ và $x+y=4$ suy ra $2x=\dfrac{13}{3}$, $x=\dfrac{13}{6}$; $2y=4-\dfrac13=\dfrac{11}{3}$, $y=\dfrac{11}{6}$. Đúng. Nhưng khoảng cách $NG$ phải bằng khoảng cách từ $G$ đến $d$. Kiểm tra $N\in d$: $\dfrac{13}{6}+\dfrac{11}{6}=4$. Đúng. Vectơ $GN=\left(\dfrac{13}{6}-\dfrac23;\dfrac{11}{6}-1\right)=\left(\dfrac32;\dfrac56\right)$. Tích vô hướng với $\overrightarrow{n}=(1;1)$: $\dfrac32+\dfrac56=\dfrac{14}{6}=\dfrac73\neq0$?? Phải vuông góc với $d$ tức song song $\overrightarrow n$. Nhưng $\overrightarrow{GN}$ không song song $(1;1)$ vì $\dfrac32\neq\dfrac56$. Vậy đường vuông góc sai.

Đường thẳng qua $G$ vuông góc với $d$ (nhận $\overrightarrow{n}=(1;1)$ làm vectơ chỉ phương) có phương trình: $\dfrac{x-\frac23}{1}=\dfrac{y-1}{1}$, tức $x-y=-\dfrac13$, hay $x-y+\dfrac13=0$. (Trước đó tôi ghi nhầm dấu.)

Giải hệ: $x+y=4$, $x-y=-\dfrac13$. Suy ra $2x=4-\dfrac13=\dfrac{11}{3}$, $x=\dfrac{11}{6}$; $2y=4+\dfrac13=\dfrac{13}{3}$, $y=\dfrac{13}{6}$.

Vậy $N\left(\dfrac{11}{6};\dfrac{13}{6}\right)$. Kiểm tra: $GN=\left(\dfrac{11}{6}-\dfrac23;\dfrac{13}{6}-1\right)=\left(\dfrac76;\dfrac76\right)$, song song $(1;1)$. Đúng.

$NG=\sqrt{\dfrac{49}{36}+\dfrac{49}{36}}=\dfrac{7\sqrt2}{6}\approx1{,}6499$. Khớp với khoảng cách từ $G$ đến $d$.

Vậy $P_{\min}=3NG=\dfrac{7\sqrt2}{2}\approx4{,}95$.
}

\end{bt}
